\section{决定性实验：先找可证伪见证，再比较算法}
\subsection{本轮已执行与尚未执行}
已执行的是第 3 节的一维非磁共振板精确核验、完整 PEC 目标级联以及有限带宽扫描。代码使用 NumPy/解析传输矩阵，没有训练、没有昂贵 Maxwell campaign，也没有运行新的二维或三维反演。

以下 P0/P1 是\textbf{预先定义的下一轮协议建议}，不是实验结果。其阈值是立项门槛，必须在新数据与随机种子生成前冻结；失败后不能移动阈值或增大噪声来救故事。

\subsection{六类场景如何服务不同命题}
\begin{longtable}{p{12mm}p{41mm}p{58mm}p{31mm}}
\toprule
场景 & 内容 & 主要判别机制 & 优先级\\
\midrule
A & 均匀介质墙 & 参考面、相位、背景模型与全波一致性 & 已有一维核验；二维回归\\
B & 未知数量/形状空腔 & 空气成核、局部补偿与 prior 支持 & P0 主场景的一部分\\
C & 倾斜/曲面/厚度变化 & 几何--定位补偿与固定深度失效 & P0 主场景\\
D & 未知层数与材料混合 & 表示能力与数据选择的分离 & P1 后验证\\
E & 导电夹杂/钢筋 & 屏蔽、共振、秩增长及网格误差 & 后置压力测试\\
F & 数据近等价而传播不同 & Candidate I 的有限非线性下界 & 第一优先\\
\bottomrule
\end{longtable}
不必先完成六类漂亮图片。P0 用一个 B+C 混合背景寻找 F 类见证，足以决定是否投入概率算法。E 需可靠金属模型和网格，不应成为首轮预算吞噬点。

\subsection{P0-A：解析回归，已经完成}
使用本研究包的 \texttt{analytic\_witness.py}。必须保留同参考平面、双程响应和 target-present 联合验证。通过条件包括：反射零点、目标响应相位差、联合目标位移补偿、lossless unitarity、频带扩展后的距离变化。独立 Fresnel 公式交叉核验与传输矩阵一致到约 $10^{-16}$。

\textbf{该门槛只保证实验定义没有初等错误。} 不把 P0-A PASS 直接当成未知复杂墙体研究 GO。

\subsection{P0-B：最小二维 TM 全波传播歧义见证}
\paragraph{固定物理。}
取 $z$ 不变的 $E_z$ 极化、$\mu=\mu_0$，求解
\[
(\nabla_{xy}^2+k_0^2\epsilon_r(x,y,\omega))E_z
=-i\omega\mu_0J_z.
\]
它是该对称条件下的 Maxwell 标量子问题，保留全部二维多次散射；并非三维矢量结果。

以 $\lambda_0$ 为单位，可取总未知域
$[-2,2]\times[0,3]$，天线位于前侧 $y=-0.3$、横向孔径 $3\lambda_0$。固定 4 个发射位置与 12 个接收位置，全都位于同一侧。先用 5 个频率 $f/f_0\in\{0.8,0.9,1,1.1,1.2\}$；另有事先冻结的窄带配置用于解释带宽作用，不能看到失败后才改采集。

材料先验覆盖全未知域，不给墙 mask；首轮限制 $1\le\Re\epsilon_r\le9$ 与明确的小损耗范围。逆表示可用 $32\times48$ 全材料/level-set 网格，相关长度、最小可表达特征和边界 prior 全部公开。truth 是一个倾斜且厚度变化的介质体加不规则空气区；这只是固定测试实例，算法不读取其生成类别或孔数。

前向网格须满足预先定义的物理精度，而非只报网格尺寸：至少做一次独立细化和 PML 扩展，要求校准数据误差的白化范数 $\le0.1$，任务响应误差 $\le0.1$；不满足则当前见证不合格。有限几何、频率和损耗要与微分算子及介质界面离散一致，避免 inverse crime。

\paragraph{两种观测合同。}
本实验沿用总场的计算分块：$W=m|_{D\setminus T}$、$X=m|_T$，其中 $T$ 是预先声明的查询域，不是已知墙体边界。在 joint 推断中，$T$ 内允许完整材料，不假设只有一种稀疏目标；传播计算中的背景才将 $T$ 换为参考材料。下述 calibration fixture 额外把 $T$ 实际设为参考材料，供隔离机制使用，这项额外信息不转移到原 joint 任务。
\begin{enumerate}
\item \textbf{Calibration 合同：} 测量时 $T$ 内实际为声明参考介质，比较 $F_{\rm cal}(W_1)$ 与 $F_{\rm cal}(W_2)$。这个合同用于隔离墙传播机制，不宣称现实任务可直接获得。
\item \textbf{Joint 合同：} 使用实际 $F(W,X,\gamma)$，两端允许目标和 nuisance 在相同 prior 内补偿。禁止从数据中减去真实墙或真实目标。它才决定原始 Research W 的适用性。
\end{enumerate}
每份合同必须有独立的 observation identity。源幅和噪声尺度由事先的校准规则固定，对所有被比较方法相同。

\paragraph{任务类和目标函数。}
任务域是前述计算查询域，不假设实际 joint 世界中不存在墙。作为固定的数值任务类，预先给定 12 个空间位置与有限组合法介质材料 $V_X$（例如参考介质内固定尺度圆盘的两种介电对比），并固定对比幅度。评价完整同任务材料全波响应，不能用任意放大 probe 获得巨大 $d_T$。此时 $X$ 是 $T$ 内全部材料，而非与 $W$ 任意相加的语义标签。

可使用局部危险方向提出初猜，再做有限非线性优化：
\[
\max_{W_2\in\Theta_W}
d_T(W_1,W_2)
\quad\text{s.t.}\quad
\|F_{\rm cal}(W_2)-F_{\rm cal}(W_1)\|_{\Sigma^{-1}}\le1,
\]
并限制物理材料范围、弱结构 prior 和最小/最大修改尺度。Joint 合同则直接优化
\[
\min_{X_2,\gamma_2}
\|F(W_2,X_2,\gamma_2)-F(W_1,X_1,\gamma_1)\|_{\Sigma^{-1}},
\]
同时保留真实任务差异（位置、占据事件或强目标传播）。位置/占据必须预先定义为总材料泛函 $q_T(m)$，例如给定对比阈值的占据事件或区域内对比质心；若采用命名物体位置，则需另声明对象生成先验。两个总材料世界必须确实不同，不能仅靠 $\chi_W+h,\chi_X-h$ 的标签重分配达标。$W_1,X_1$ 是声明的参考世界；该构造是可辨识性 adversarial test，不是在线重建性能。

对照为相同预算的普通 Full-J/adjoint、材料匹配随机方向加共同 corrector、multi-start joint FWI。不得把人为阻止普通方法进展的角度约束放入对照。各方法可使用相同 factorization 与 warm-start 权利。

\paragraph{预算与先后顺序。}
首轮只做两个固定背景种子和两种冻结采集配置，先 calibration，再对通过的见证检验 joint。建议整个 P0-B 硬上限为 8,000 complex RHS、2,000 个材料--频率分解和 2 CPU 小时，任一先到即止；40\% 用搜索，30\% 用完整端点/任务验证，20\% 用独立网格，10\% 留失败/重查。具体后端若无法在此预算内完成，缩减场景数量而非把失败费用排除。

这些数值是起始预算，不是理论复杂度保证。只在执行前根据一次公开的 solver benchmark 调整；不能在看到科学结果后重新制定上限。

\paragraph{GO / NO-GO。}
\begin{longtable}{p{24mm}p{124mm}}
\toprule
门槛 & 判定\\
\midrule
机制 GO & 至少一个有限尺寸、非磁、被动、不规则二维 pair：数据均值距离 $d\le1$，事先声明目标类的完整响应有 $d_T\ge3$，并通过独立细化与 nuisance 合法性；报告其是否仅 calibration 成立。\\
原问题 GO & 在 target-present joint 合同下仍有 $d\le1$，且目标任务差异达到事先物理尺度，例如位置差 $\ge0.1\lambda_0$ 或声明占据事件改变；不能用纯墙距离替代。\\
稳健性检查 & 在预定的微小材料/位置扰动下继续检查距离，披露有效范围。孤立精细调谐反例可支持存在性，不能自动支持常见场景算法价值。\\
机制 NO-GO & 所有有限非线性正确器都消去传播差异，或只有错误边界/粗网格/不合法 gain 才满足；停止算法立项。\\
范围 REDIRECT & 仅单频或极窄带成立，实际冻结宽带 joint 配置均不成立：缩为测量充分性/反例边界论文，不继续“通用复杂墙递进”叙事。\\
搜索不确定 & 固定预算没找到：报告 best-found 和失败状态，不能证明全局不存在；但可以据资源决策停止该算法分支。\\
\bottomrule
\end{longtable}

$d\le1$ 对等 prior 固定 Gaussian 两世界意味着最优检验错误率至少 $\Phi_{\rm N}(-1/2)\approx0.309$。这是人为选择的强歧义门槛，不是“95\% noise ball”。$\|y-F\|$ 接近 $\sqrt{\dim y}$ 的两个拟合可有很大相互距离，Research T 的历史结果已说明必须分开。\cite{R_T}

数值容差必须进入门槛。若各端有严格的响应误差上界，应要求
\[
d_{\rm num}+\epsilon_{F_1}+\epsilon_{F_2}\le1,\qquad
d_{T,\rm num}-\epsilon_{H_1}-\epsilon_{H_2}\ge3.
\]
仅观察网格差 $\le0.1$ 不能保护恰好卡在阈值上的见证，也不是连续 PDE 误差的严格上界。若没有可靠误差证书，结论须限定为最细声明的离散 Maxwell 模型，并把网格/PML 检查列为数值一致性证据。

\subsection{P1：区分多解收益与推进收益}
只有 P0-B 支持实际 joint 歧义，才进入 P1。首先用两至数个固定墙分支加少量连续目标/校准参数建立可独立积分的 reference。它检验概率逻辑和压缩界，不声称任意未知拓扑 posterior。原 prior 的分支权重必须明确，后验权重通过 evidence 积分得到，不能用最小残差替代。

\begin{center}
\begin{tabular}{p{29mm}p{55mm}p{62mm}}
\toprule
 & 同时联合计算 & 自适应分块/递进计算\\
\midrule
单点近似 & 正则 joint FWI/MAP & 单点交替更新传播与材料，允许回访\\
多假设 & 直接 joint SMC/MH/参考积分 & 同一 posterior 上的任务风险传播缓存和调度\\
\bottomrule
\end{tabular}
\end{center}
这四格将多假设因素与推进因素分开。Oracle 真墙 Green、错误均匀背景和 free-space 成像另外列为诊断上下参照；它们没有相同可用信息，不能与主方法算“公平胜出”。

主多假设对照使用同一 prior、同一完整 likelihood、相同任务和 proposal 库；优化/提案顺序可以不同，但总费用相同。若多假设已经改善，而递进调度没有进一步改善，应只保留多解传播贡献。

\paragraph{低成本 pilot。}
先 4 个开发背景定参数，再冻结 8 个背景，每背景两份噪声；所有指标以背景为聚合单位。每个方法先锁定相同 full-evaluation-equivalent 上限（建议 256 个多频多源完整前向等价，且同时限实际 wall time），参考积分的费用单列但不免费消失。若短链无法达到参考误差门槛，P1 应判未解决，不用粒子 ESS 冒称校准。

在低维 reference 上，核心 GO 是：重要事件质量及目标 CDF 在固定检查点误差 $\le0.02$；传播压缩保持同一误差目标且实际费用比直接 full 多假设低至少 20\%，或在同费用下显著降低声明任务损失。阈值须在执行前冻结。8 个背景的 coverage 仅作 pilot 描述，不能据少量命中给现实覆盖保证。

\paragraph{必须报告的指标。}
\begin{itemize}
\item 在线决策的目标位置/形状/材料误差；决策为 posterior mean、medoid 或明确 Bayes action，不用 truth-selected Best-of-K。
\item 目标事件概率误差、proper score、校准曲线、可信区间 coverage 及宽度；有独立 reference 才谈概率准确度。
\item 完整目标传播预测误差、四块签名误差与共振裕度；数值误差和环境不确定性分列。
\item 重要模式发现和遗漏、错误固定率、回退次数；无法估计未覆盖质量时填 unknown。
\item 原始全 Maxwell 数据一致性、留出测量误差、所有失败状态。
\item 总时间、峰值内存、分解/RHS、basis/setup、全部拒绝与 fallback；同 rank 不替代同费用。
\end{itemize}
Best-of-K 可以作为明确离线容量诊断，但不能成为在线性能主指标。A30、Research A 与 P2 已表明这个区别足以改变科研结论。\cite{R_A30,R_A,R_P2}

\subsection{独立检验压缩证书与 A17 价值}
为四块签名构造若干受控误差，分别比较仅透射、Born 核、完整四块、普通 goal-ROM 与 A17 条件 enrichment。包含远离共振和接近共振两类；必须允许证书在后者失效。

用实际 full target response 验证上界是否保守、是否错误声称安全，以及界的求值是否比 direct solve 更贵。若只有完整 reference 已知时才能决定 fallback，称为 offline audit；不得把它包装为在线控制。P2 已有这一失败先例。\cite{R_P2}

A17 分支只有在完整总费用对强 direct/Krylov/goal-ROM 存在任务误差--成本 Pareto 改善时继续。若相同空间数值相同、额外选空间却更贵，即停止其优势主张；这不推翻 A17 的冻结机制结果。

\section{现实测量的最低条件与三维推广}
\subsection{不存在脱离任务的统一最低频段}
需要的带宽、孔径和 SNR 应由目标任务距离、材料损耗和 nuisance 共同确定。均匀已知介质中的近似距离分辨率
\[
\Delta z\approx\frac{c}{2B\sqrt{\epsilon_r}}
\]
只能作为量级参考；未知弯曲路径、色散、多次散射和介电误差会改变定位。横向分辨率也依赖角覆盖与距离，不能从中心波长给出普遍数值。

实际最低合同应包括：稳定且可建模的发射谱；复相位可追踪的多静态数据；噪声/相关性估计；阵列位置与时间零点的不确定范围；至少能使已知环境下的目标任务可测的 SNR。若连 Oracle 真墙条件下都无法区分目标任务，研究环境学习没有可挽回的观测信息。

墙内衰减可用路径积分作局部解释：单程幅度约 $e^{-\int\alpha ds}$，往返目标响应约 $e^{-2\int\alpha ds}$；相应均值导数平方的信息量可按 $e^{-4\int\alpha ds}$ 衰减。这是传播路径近似，复杂共振或旁路需 full Maxwell。更高频率可能增加分辨率，也增加衰减与模型敏感性，不能预设总是更好。

\subsection{单侧、多静态、极化与门控}
主验证保持所有天线在同侧。多静态与多极化可改变可达/可观测子空间，但若硬件相位或极化通道缺校准，也增加 nuisance。比较须保留真实通道协方差和重复采集的总能量/时间预算。

早时门控只有在目标最早响应确实晚于保留窗口、且有限带宽泄漏有界时，才可近似得到墙主导数据。未知外表面、侧向绕射、薄层共振和墙--目标耦合会破坏分离。门控丢弃数据通常不会增加信息；它可以改善数值动态范围，但要在 likelihood 中计入剩余模型误差和噪声相关性。

不存在准确 CAD/LiDAR 时，传感器位置本身仍需相对坐标参考。可以把配准误差作为 nuisance；不能既不测量又假定完全准确。额外 LiDAR、背面接收或已知反射体只能作为明确新增信息实验，不能偷偷成为原问题前提。\cite{L_Lidar}

\subsection{从二维 TM 到三维全矢量}
四块签名、Schur/resolvent、推送后验和统计下界可在相应 Hilbert 空间/离散向量系统中保持形式。三维需要重新解决：
\begin{enumerate}
\item $H(\mathrm{curl})$ 相容离散、散度约束、矢量 Green 自项及边界条件；不能逐像素独立标量拟合后称矢量一致。
\item 多源多极化共用同一张量/标量材料场与色散模型；current proposal 必须共同可实现。
\item 对金属的皮肤深度、表面电流与连通网络；高导电极限不能直接沿用粗二维参数。
\item 更大的端口秩、factorization 与预条件成本，非正规/共振对误差估计的影响。
\item 开放有限孔径的三维边缘绕射和出平面泄漏，它们可能消除二维歧义，也可能创造新歧义。
\end{enumerate}
因此二维 PASS 只支持机制继续验证。最终 TAP 若主张一般建筑穿墙，应至少有一组有限三维矢量、同侧采集的物理验证；没有则题名和结论应限定二维。

